% !TEX program = pdflatex
% ------------------------------------------------------------------
% The exact DNF complexity of the inner product function
% LaTeX version of the Markdown master notes/dnf_min_ip_arxiv_en.md
% (NOTE-001, 2026-07-23). Author: Tanguy Le Mer.
% Gate before submission: independent expert proofreading (see the
% French master note, reserve journal §6).
% Citations are kept inline, faithful to the master; exact paper
% titles to be completed by the author at proofreading time
% (no bibliographic detail invented).
% ------------------------------------------------------------------
\documentclass[11pt]{article}

\usepackage[T1]{fontenc}
\usepackage[utf8]{inputenc}
\usepackage{amsmath,amssymb,amsthm}
\usepackage{booktabs}
\usepackage[margin=1.1in]{geometry}
\usepackage[colorlinks=true, linkcolor=blue!50!black,
            urlcolor=blue!50!black]{hyperref}

\newtheorem{theorem}{Theorem}
\newtheorem*{remark}{Remark}

\newcommand{\IP}{\mathrm{IP}}
\newcommand{\DNFmin}{\mathrm{DNF}_{\min}}
\newcommand{\CNFmin}{\mathrm{CNF}_{\min}}
\newcommand{\primes}{\mathrm{primes}}
\newcommand{\WL}{W_{L_1}}

\title{The exact DNF complexity of the inner product function}
\author{Tanguy Le Mer}
\date{}

\begin{document}
\maketitle

\begin{abstract}
Let $\IP_{2n}(x,y)=\sum_{i=1}^{n} x_i y_i \bmod 2$ be the inner
product function on $2n$ variables. We determine its exact DNF
complexity: the minimum number of terms in a DNF (OR of ANDs of
literals, unbounded fan-in) computing $\IP_{2n}$ is exactly
$(3^n-1)/2$, and the minimum DNF is unique --- $\IP_{2n}$ has exactly
$(3^n-1)/2$ prime implicants, each of which is essential. Dually,
$\CNFmin(\IP_{2n})=(3^n+1)/2$. We further prove a correspondence
theorem: the minimum $L_1$ weight of a real polynomial threshold
representation of $\IP_{2n}$ (monomial basis over $\{0,1\}$) equals
$(3^n-1)/2$ as well, with matching upper bound (the exact multilinear
polynomial) and lower bound (an explicit diagonal dual certificate all
of whose subcube moments are $\pm 1$). In terms of the total number of
variables $N=2n$, the DNF size is $(3^{N/2}-1)/2 \approx 2^{0.7925N}$
--- strictly between polynomial size and the generic $2^{N-1}$ of
parity. After an extended literature search (detailed in
Section~\ref{sec:related}) we believe the exact constant is not
recorded; the result complements the active line of work on the
depth-3 complexity of inner product by settling the depth-2 base case
exactly. All statements are verified by exhaustive computation and
exact-arithmetic certificates at small sizes.
\end{abstract}

\section{Statement}

\begin{theorem}[DNF]\label{thm:dnf}
$\DNFmin(\IP_{2n}) = (3^n-1)/2$. Moreover $\IP_{2n}$ has exactly
$a_n = (3^n-1)/2$ prime implicants, each essential (each admits a
witness point covered by it alone); the minimum DNF is therefore
unique and equals the disjunction of all prime implicants. The
companion count for the complement is
$\lvert\primes(\lnot\IP_{2n})\rvert = b_n = (3^n+1)/2$.
\end{theorem}

\begin{theorem}[CNF]\label{thm:cnf}
$\CNFmin(\IP_{2n}) = (3^n+1)/2$.
\end{theorem}

\begin{theorem}[weight-PTF correspondence]\label{thm:ptf}
For every $n \ge 1$, the minimum $L_1$ coefficient weight of a real
polynomial $p$ (monomial basis over $\{0,1\}^{2n}$, unbounded degree)
with $p \ge 1$ on $\IP^{-1}(1)$ and $p \le 0$ on $\IP^{-1}(0)$ is
exactly $(3^n-1)/2 = \DNFmin(\IP_{2n})$.
\end{theorem}

\section{Proof of Theorem~\ref{thm:dnf}}

The engine is the disjoint XOR bi-decomposition
\begin{equation}\label{eq:bidec}
\IP_{2(n+1)}(x, y, x_{n+1}, y_{n+1})
  = \IP_{2n}(x, y) \oplus (x_{n+1} \wedge y_{n+1}).
\end{equation}

\paragraph{(A) Classification of prime implicants.}
Every prime implicant of $\IP_{2(n+1)}$ is of exactly one of three
forms:
\begin{itemize}
\item \textbf{A1}: $x_{n+1} \wedge y_{n+1} \wedge M_0$, with $M_0$ a
  prime implicant of $\lnot\IP_{2n}$;
\item \textbf{A2a}: $\lnot x_{n+1} \wedge M_0$, with $M_0$ a prime
  implicant of $\IP_{2n}$;
\item \textbf{A2b}: $\lnot y_{n+1} \wedge M_0$, with $M_0$ a prime
  implicant of $\IP_{2n}$.
\end{itemize}
Indeed, a term leaving the pair $(x_{n+1}, y_{n+1})$ free would force
$\IP_{2n} = \lnot\IP_{2n}$ on its projection; and
$\lnot x_{n+1} \wedge \lnot y_{n+1} \wedge M_0$ is never prime (it
weakens to A2a or A2b).

\paragraph{(B) Recurrence and closed form.}
Writing $a_n = \lvert\primes(\IP_{2n})\rvert$ and
$b_n = \lvert\primes(\lnot\IP_{2n})\rvert$, the classification gives
\[
a_{n+1} = 2a_n + b_n, \qquad b_{n+1} = a_n + 2b_n,
\qquad (a_1, b_1) = (1, 2),
\]
whence $a_n + b_n = 3^n$ and $b_n - a_n = 1$ by induction, so
$a_n = (3^n-1)/2$ and $b_n = (3^n+1)/2$.

\paragraph{(C) Essentiality (lower bound).}
For each prime implicant $M$ we construct a point covered by $M$ and
by no other prime implicant, by induction. If
$M = \lnot x_{n+1} \wedge M_0$ (class A2a), take the inductive witness
$x_0$ of $M_0$ in $\IP_{2n}$ and append
$(x_{n+1}, y_{n+1}) = (0, 1)$; the fixed pair excludes all A1 terms
and all A2b terms with the mirrored coordinate, and the inductive
property excludes competing A2a terms. Symmetrically for A2b with
$(1, 0)$, and for A1 with $(1, 1)$ using the witness of $M_0$ in
$\lnot\IP_{2n}$. Every DNF for $\IP_{2(n+1)}$ must therefore contain
every prime implicant, so $\DNFmin = a_n$. \qed

\section{Proof of Theorem~\ref{thm:cnf}}

By De Morgan, $\CNFmin(f) = \DNFmin(\lnot f)$. The key observation is
that the XOR bi-decomposition~\eqref{eq:bidec} is invariant under
complementation:
\[
\lnot\IP_{2(n+1)} = \lnot\IP_{2n} \oplus (x_{n+1} \wedge y_{n+1}),
\]
i.e.\ the complement has exactly the same shape
$g \oplus (x \wedge y)$ with base $\lnot\IP_{2n}$. The entire
apparatus of the previous section transfers verbatim with
$\IP \leftrightarrow \lnot\IP$ exchanged: the classification yields
the second half of the recurrence ($b_{n+1} = a_n + 2b_n$), and the
witness construction yields essentiality of every prime implicant of
$\lnot\IP_{2n}$ (base case $n = 1$:
$\lnot\IP_2 = \lnot x_1 \vee \lnot y_1$, two primes, witnesses
$(0,1)$ and $(1,0)$). Hence
$\DNFmin(\lnot\IP_{2n}) = b_n = (3^n+1)/2$. \qed

\section{Proof of Theorem~\ref{thm:ptf}}

\paragraph{Upper bound.}
The exact multilinear polynomial
$p^{*} = \bigl(1 - \prod_i (1 - 2 x_i y_i)\bigr)/2$ has coefficients
$(-1)^{\lvert S\rvert + 1}\, 2^{\lvert S\rvert - 1}$ on the paired
monomials $\prod_{i \in S} x_i y_i$, so its $L_1$ weight is
$\sum_{s \ge 1} \binom{n}{s} 2^{s-1} = (3^n-1)/2$, and $p^{*}$ takes
values exactly in $\{0, 1\}$, hence is a valid threshold
representation ($p = p^{*}$ works since $p^{*} \in \{0,1\}$).

\paragraph{Lower bound (diagonal dual certificate).}
Consider the signed measure supported on the diagonal $x = y$:
\[
w(x, x) = (-1)^{\lvert x\rvert + 1}\, 2^{\,n - \lvert x\rvert}
\quad \text{for } \lvert x\rvert \ge 1,
\qquad w(0, 0) = -(2^n - 2),
\]
and $w = 0$ off the diagonal. A one-line computation shows every
subcube moment of $w$ equals $\pm 1$ (the inner sum telescopes to
$(-1)^{s+1}\, 2^{\,n-s-r}\, (2-1)^r$ with $n - s - r = 0$), and its
positive mass is
$\sum_{a \text{ odd}} \binom{n}{a} 2^{\,n-a} = (3^n-1)/2$. Weak LP
duality then gives $\WL \ge (3^n-1)/2$ over the reals. \qed

\medskip
The certificate lives on the diagonal $x = y$ --- where $\IP$
restricts to parity. Note the coincidence $\DNFmin = \WL$: a priori
the two quantities are incomparable (neither bound implies the
other); $\IP$ realizes both at the same value.

\paragraph{A transversal constant.}
Three depth-2-type quantities for $\IP$ land on the same scale
$3^{n}$:

\begin{center}
\begin{tabular}{lll}
\toprule
Model & Quantity & Value \\
\midrule
DNF & $\DNFmin(\IP_{2n})$ & $(3^n-1)/2$ (Thm~\ref{thm:dnf}) \\
CNF & $\CNFmin(\IP_{2n})$ & $(3^n+1)/2$ (Thm~\ref{thm:cnf}) \\
Consistent 2-CNF & max satisfying assignments & $3^n$ (tight;
  Frankl--Gryaznov--Talebanfard, ITCS 2022) \\
Real PTF weight & $\min \WL(\IP_{2n})$ & $(3^n-1)/2$
  (Thm~\ref{thm:ptf}) \\
\bottomrule
\end{tabular}
\end{center}

\section{Related work and novelty status}\label{sec:related}

The active line on the \textbf{depth-3} complexity of $\IP$
(Golovnev--Kulikov--Williams, ITCS 2021; Frankl--Gryaznov--%
Talebanfard, ITCS 2022; G\"o\"os--Guan--Mosnoi, MFCS 2023 and
Information \& Computation 2024,
$\Sigma_2^3 \in [2^{0.847n},\, 2^{0.965n}]$; Amano, ISAAC 2023;
Gurumukhani et al., 2026) never states the depth-2 case; this note
supplies its exact base point. Standard covering bounds give only
$\sim 2^n - 1$ (the essentiality of \emph{all} primes is what pushes
to $3^n/2$). The switching-theory literature was checked at the
source: Sasao--Butler (ASP-DAC 2001) handle AND/OR-orthodox
bi-decompositions --- the XOR case of $\IP$ is outside their calculus
(verified on the full text); the Butler--Sasao cryptographic tutorial
treats $\IP$ as a bent function from the ANF angle, with no SOP size.
Compendia (Jukna 2012; Crama--Hammer 2011) state the prime-implicant
framework but not this count. A targeted search of the Russian school
(Yablonskii and collaborators, \emph{Problemy Kibernetiki}), in
Russian and English, found no trace of the function or the constant;
non-digitized volumes remain the sole residual reserve. We therefore
present the exact constant as, to the best of our knowledge,
unrecorded --- and would welcome any earlier reference.

\section{Computer verification}

All claims are machine-verified in the accompanying repository:
$\DNFmin = 1, 4, 13, 40, 121$ for $n = 1..5$ by exhaustive prime
enumeration plus exact set-cover ILP (CBC); prime counts and
essentiality on both $\IP$ and $\lnot\IP$ for $n = 1..4$ by brute
force; the dual certificate of Theorem~\ref{thm:ptf} in exact
rational arithmetic for $n = 1..14$ plus brute-force cross-check for
$n \le 4$; the LP optimum $40/121$ at $n = 4, 5$ matched by the
symmetrized dual to $10^{-8}$.

\section{Disclosure}

This note results from an AI-assisted exploration (Claude,
Anthropic), with the adversarial-verification methodology of the
repository \texttt{pvsnp-exploration}: every claim above is either
proved in full in this note or backed by an executable
exact-arithmetic check; the proofs were human-read by the author. AI
assistance is disclosed per arXiv policy.

\end{document}
